% =====================================================================
% 00-Foundations.tex
% Geometric Hallucination Taxonomy — Mathematical Spine
% Borealis S. Hedling — 2026
% =====================================================================

\section*{00. Foundations}

This document establishes the geometric foundations used throughout the
mathematical spine of the Geometric Hallucination Taxonomy (GHT). All
operators defined in subsequent modules rely on the structures introduced
here: smooth manifolds, Riemannian metrics, Levi–Civita connections,
tangent bundles, and latent trajectories. The presentation follows
standard differential–geometric conventions and is compatible with
mechanized verification in Lean (Appendix B).

% ---------------------------------------------------------------------
\subsection*{0.1 Smooth Manifold}

Let $M$ be a smooth manifold of dimension $n$, modeled on $\mathbb{R}^n$.
We denote its tangent bundle by $TM$ and write $T_p M$ for the tangent
space at a point $p \in M$.

Charts are written as $(U, \varphi)$ with $\varphi : U \to \mathbb{R}^n$
a diffeomorphism onto its image. We assume $M$ is Hausdorff and second
countable.

% ---------------------------------------------------------------------
\subsection*{0.2 Riemannian Metric}

A Riemannian metric on $M$ is a smooth assignment
\[
g : p \mapsto g_p,
\]
where each $g_p$ is a positive–definite inner product on $T_p M$.
In local coordinates $(x^1, \dots, x^n)$, the metric is represented by
a symmetric matrix $(g_{ij})$ with
\[
g_{ij}(p) = g_p\!\left( \frac{\partial}{\partial x^i},
                       \frac{\partial}{\partial x^j} \right).
\]

We denote the inverse metric by $(g^{ij})$.

% ---------------------------------------------------------------------
\subsection*{0.3 Levi–Civita Connection}

The Levi–Civita connection $\nabla$ is the unique torsion–free connection
compatible with the metric:
\[
\nabla g = 0, \qquad T(X,Y) = \nabla_X Y - \nabla_Y X - [X,Y] = 0.
\]

In coordinates, the Christoffel symbols are
\[
\Gamma^i_{\;jk}
= \frac{1}{2} g^{im}
\left(
\partial_j g_{km}
+ \partial_k g_{jm}
- \partial_m g_{jk}
\right).
\]

Covariant differentiation of a vector field $Y = Y^i \partial_i$ is
\[
\nabla_X Y
= X^j \left( \partial_j Y^i + \Gamma^i_{\;jk} Y^k \right) \partial_i.
\]

% ---------------------------------------------------------------------
\subsection*{0.4 Tangent Vectors and Norms}

For $v \in T_p M$, the norm induced by the metric is
\[
\|v\|_g = \sqrt{g_p(v,v)}.
\]

Inner products are written $\langle u, v \rangle_g$.

% ---------------------------------------------------------------------
\subsection*{0.5 Latent Trajectories}

A latent trajectory is a smooth curve
\[
x : I \to M,
\]
where $I \subseteq \mathbb{R}$ is an interval representing inference
steps. In discrete settings (as used in the Python toolkit), we write
\[
x_{t+1} = \Phi(x_t),
\]
where $\Phi : M \to M$ is a smooth update rule.

The velocity vector is
\[
\dot{x}(t) = \frac{d}{dt} x(t) \in T_{x(t)} M.
\]

% ---------------------------------------------------------------------
\subsection*{0.6 Geodesics}

A geodesic $\gamma$ satisfies
\[
\nabla_{\dot{\gamma}} \dot{\gamma} = 0.
\]

In coordinates:
\[
\ddot{\gamma}^i + \Gamma^i_{\;jk} \dot{\gamma}^j \dot{\gamma}^k = 0.
\]

Geodesics are used throughout the taxonomy to measure curvature–driven
instability.

% ---------------------------------------------------------------------
\subsection*{0.7 Parallel Transport}

Given a curve $\gamma : [0,1] \to M$, parallel transport of a vector
field $V(t)$ along $\gamma$ satisfies
\[
\nabla_{\dot{\gamma}(t)} V(t) = 0.
\]

The solution defines a linear map
\[
P_\gamma : T_{\gamma(0)} M \to T_{\gamma(1)} M.
\]

Parallel transport is the foundation for holonomy drift analysis.

% ---------------------------------------------------------------------
\subsection*{0.8 Metric Rank and Degeneracy}

The metric tensor $g$ may lose rank in unstable regions. We define the
rank at $p$ as
\[
\tau(p) = \operatorname{rank}(g_p).
\]

A collapse surface is a set
\[
\Sigma = \{ p \in M : \det(g_p) = 0 \}.
\]

These structures underpin topological instability analysis.

% ---------------------------------------------------------------------
\subsection*{0.9 Fusion–Field Preliminaries}

Given $x, y \in T_p M$, a fusion–field is a smooth map
\[
F : TM \times TM \to TM
\]
that blends representations. Nonlinear interaction terms are denoted
$\sigma(x,y)$ and must satisfy smoothness and bilinearity constraints
specified in Module 05.

% ---------------------------------------------------------------------
\subsection*{0.10 Stability Thresholds}

Each operator in the taxonomy uses a stability threshold:
\[
\kappa_{\mathrm{thr}}, \quad
\sigma_{\mathrm{thr}}, \quad
h_{\mathrm{thr}}, \quad
\tau_{\mathrm{thr}}, \quad
f_{\mathrm{thr}}.
\]

These thresholds define instability regimes and are treated as fixed
positive constants throughout the spine.

% ---------------------------------------------------------------------
\subsection*{0.11 Notational Conventions}

\begin{itemize}
\item $X,Y,Z$ denote vector fields.
\item $u,v$ denote tangent vectors.
\item $g_{ij}$ and $g^{ij}$ denote metric components.
\item $\Gamma^i_{\;jk}$ denote Christoffel symbols.
\item $R^i_{\;jkl}$ denotes curvature tensor components.
\item $\lambda(t)$ denotes eigenvalues of latent representations.
\item $P_\gamma$ denotes parallel transport.
\item $H = P_\gamma - I$ denotes holonomy.
\item $\tau$ denotes metric rank.
\item $F(x,y)$ denotes fusion–field output.
\end{itemize}

These conventions remain consistent across all modules.

% =====================================================================
% End of 00-Foundations.tex
% =====================================================================
